% MyEspresso showcase: every environment and element the theme colours.
%
% One PDF per colour scheme, the default (red) when nothing is predefined:
%   latexmk -lualatex -jobname=showcase-blue \
%     -usepretex='\def\MEscheme{myespressoblue}\def\MElabel{blue}' showcase.tex
%
% Side by side, each slide in red, blue, green and black (generated, not tracked):
%   args=(); for i in {1..15}; do for s in red blue green black; do
%     args+=(showcase-$s.pdf $i); done; done
%   pdfjam "${args[@]}" --nup 2x2 --papersize '{320mm,180mm}' --frame true \
%     --outfile showcase-compare.pdf
\providecommand\MEscheme{}
\providecommand\MElabel{red}
\documentclass[aspectratio=169]{beamer}

\edef\MEoptions{\MEscheme,delaunay,randomseed=1}
\expandafter\usetheme\expandafter[\MEoptions]{MyEspresso}

\newtheorem{proposition}[theorem]{Proposition}

% A swatch row: colour, role, where the theme uses it. The hex is read back
% from xcolor in the last cell, since each cell is its own group.
\newcommand\swatch[3]{%
  \textcolor{#1}{\rule{2.4em}{1.1em}} & #2 & #3 &
  \extractcolorspecs{#1}{\swm}{\sws}\convertcolorspec{\swm}{\sws}{HTML}\swh
  \texttt{\#\swh}\\}

\title[MyEspresso \MElabel]{MyEspresso Showcase}
\subtitle{Colour scheme: \MElabel}
\author{A.~Author}
\institute{Department of Examples}
\mail{author@example.org}
\date{October 2026}

\begin{document}

\begin{frame}[plain]
  \titlepage
\end{frame}

\section{Colours and Text}
\sepframe

\begin{frame}[label=palette]{Role colours}
  \footnotesize
  \begin{tabular}{@{}lllr@{}}
    & \textbf{role} & \textbf{used by} & \\
    \swatch{structure}{structure}{titles, block titles, theorem family, covers}
    \swatch{structure!60}{structure!60}{list markers, links}
    \swatch{alerted text.fg}{theme alert}{\texttt{\textbackslash alert}, alertblock, \texttt{\textbackslash boxalert}, citations}
    \swatch{alert}{xcolor \texttt{alert}}{beamer's fixed red, \texttt{\textbackslash textcolor\{alert\}}}
    \swatch{MyEspresso@definition}{definition}{definition titles (\texttt{mydefinitioncolor})}
    \swatch{MyEspresso@example}{example}{example, exampleblock (\texttt{mypurple})}
    \swatch{MyEspresso@info}{info}{information (\texttt{mycalloutcolor})}
    \swatch{MyEspressoRemarkColor}{remark}{remark box}
  \end{tabular}
\end{frame}

\begin{frame}{Inline highlights and links}
  \begin{itemize}
    \item \texttt{\textbackslash alert}: an \alert{alerted phrase}
    \item \texttt{\textbackslash boxalert}: a \boxalert{boxed alert}
    \item \texttt{\textbackslash textcolor\{alert\}}: a \textcolor{alert}{red highlight}
    \item \texttt{\textbackslash structure}: a \structure{structure phrase}
    \item \texttt{\textbackslash emph}: an \emph{emphasised phrase}
    \item A link to the \hyperlink{palette}{role colours}, a URL \url{https://example.org}, and a citation \cite{Halmos1974}
  \end{itemize}
\end{frame}

\begin{frame}{Lists}
  \begin{columns}[T]
    \begin{column}{0.48\textwidth}
      \begin{itemize}
        \item First item
        \begin{itemize}
          \item Nested item
        \end{itemize}
        \item Second item
      \end{itemize}
      \begin{enumerate}
        \item First step
        \begin{enumerate}
          \item Nested step
        \end{enumerate}
        \item Second step
      \end{enumerate}
    \end{column}
    \begin{column}{0.48\textwidth}
      \begin{description}
        \item[Term] A description.
        \item[Longer term] Another description.
      \end{description}
    \end{column}
  \end{columns}
\end{frame}

\begin{frame}{Frame sections and buttons}
  \framesection{A frame section}
  Text under an in-frame heading.

  \medskip
  \beamerbutton{Button}\quad
  \hyperlink{palette}{\beamergotobutton{Go to}}\quad
  \beamerreturnbutton{Return}\quad
  \beamerskipbutton{Skip}
\end{frame}

\section{Theorem-like Environments}
\sepframe

\begin{frame}{The theorem family}
  \begin{columns}[T]
    \begin{column}{0.48\textwidth}
      \begin{theorem}[Main result]
        The headline statement.
      \end{theorem}
      \begin{proposition}[Supporting]
        A statement declared by the document.
      \end{proposition}
    \end{column}
    \begin{column}{0.48\textwidth}
      \begin{lemma}[Technical]
        A step on the way.
      \end{lemma}
      \begin{corollary}
        A consequence.
      \end{corollary}
      \begin{fact}
        A known fact.
      \end{fact}
    \end{column}
  \end{columns}
\end{frame}

\begin{frame}{Definitions, examples and proofs}
  \begin{columns}[T]
    \begin{column}{0.48\textwidth}
      \begin{definition}[Admitted pair]
        A named object.
      \end{definition}
      \begin{definitions}
        Several named objects.
      \end{definitions}
      \begin{proof}
        A short argument.
      \end{proof}
    \end{column}
    \begin{column}{0.48\textwidth}
      \begin{example}[Torus]
        \begin{itemize}
          \item A list inside an example
        \end{itemize}
      \end{example}
      \begin{examples}
        Several instances.
      \end{examples}
    \end{column}
  \end{columns}
\end{frame}

\section{Blocks and Asides}
\sepframe

\begin{frame}{Blocks}
  \begin{columns}[T]
    \begin{column}{0.31\textwidth}
      \begin{block}{block}
        \begin{itemize}
          \item Plain block
        \end{itemize}
      \end{block}
    \end{column}
    \begin{column}{0.31\textwidth}
      \begin{alertblock}{alertblock}
        \begin{itemize}
          \item Alerted block
        \end{itemize}
      \end{alertblock}
    \end{column}
    \begin{column}{0.31\textwidth}
      \begin{exampleblock}{exampleblock}
        \begin{itemize}
          \item Example block
        \end{itemize}
      \end{exampleblock}
    \end{column}
  \end{columns}
\end{frame}

\begin{frame}{Asides}
  \begin{columns}[T]
    \begin{column}{0.48\textwidth}
      \begin{information}
        An informational aside.
      \end{information}
      \begin{remark}[Scope]
        A remark box.
      \end{remark}
    \end{column}
    \begin{column}{0.48\textwidth}
      \begin{callout}{e\_class callout}
        The e\_class note box, framed in \texttt{mycalloutcolor}.
      \end{callout}
    \end{column}
  \end{columns}
\end{frame}

\begin{frame}{Quotation and abstract}
  \begin{quotation}[Halmos]
    Less is more.
  \end{quotation}
  \begin{abstract}
    An abstract set between two rules.
  \end{abstract}
\end{frame}

\begin{frame}{References}
  \begin{thebibliography}{9}
    \bibitem{Halmos1974} P.~R.~Halmos. How to talk mathematics. \emph{Notices AMS} 21 (1974) 155--158.
  \end{thebibliography}
\end{frame}

\thanksframe{Thank you}

\end{document}
